令 $b_k=\|S^{(k)}f_k\|_H$、$h=(\sum_k b_k^2)^{1/2}$ 和 $W=(\sum_k |f_k|^2)^{1/2}$。令 $M$ 为对立方体取平均的 Hardy--Littlewood 极大算子，并令 $L=\lceil(N+2)\log_2(2/\delta)\rceil$。对每个二进立方体 $Q$，边长为 $2^{-j}$，中心为 $y$，比较 $h$ 与 $c_Q=(\sum_{k\le j-L} b_k(y)^2)^{1/2}$。对 $x,z\in Q$ 和 $k\le j-L$，导数界与中值不等式给出
\[
 \|S^{(k)}f_k(x)-S^{(k)}f_k(y)\|_H
 \lesssim C_*\delta^{-N}2^{k-j}MW(z).
\]
事实上 $2^k|x-z|\lesssim1$，而沿从 $y$ 到 $x$ 的线段平移所得的核，都可以用 $2^{nk}(1+2^k|z-\cdot|)^{-N}$ 的固定常数倍控制。将这个径向递减核与 $W$ 卷积，所得至多为 $MW(z)$ 的常数倍。各差的平方和开方由 $CC_*\delta MW(z)$ 控制，因为 $\delta^{-N}2^{-L}\lesssim\delta$。

对于 $k\ge j+L$，将 $f_k$ 分成 $5Q$ 上的部分及其补集上的部分，其中 $5Q$ 是 $Q$ 的同心伸缩。伸缩保持 $L^2$ 算子范数，因此对 $z\in Q$，
\[
 \frac1{|Q|}\int_Q\sum_{k\ge j+L}
 \|S^{(k)}(f_k\mathbf1_{5Q})(x)\|_H^2\,dx
 \lesssim B_2^2 M(W^2)(z).
\]
对于 $x,z\in Q$，在 $5Q$ 逐次加倍得到的壳层上累加核估计，得到
\[
 \|S^{(k)}(f_k\mathbf1_{(5Q)^c})(x)\|_H
 \lesssim C_*\delta^{-N}2^{-(N-n)(k-j)} MW(z).
\]
对边长与 $2^{a-j}$ 可比的立方体，$x$ 到相应外部壳层的距离至少为 $2^{a-j}$ 的常数倍。$|f_k|$ 在该立方体上的积分可由其体积乘以 $MW(z)$ 控制。因此对 $a$ 的余下求和收敛，因为 $N>n$。再对 $k\ge j+L$ 作平方求和开方，所得由 $CC_*\delta MW(z)$ 控制。

个数少于 $2L$ 的剩余指标 $k$，在这一有限集合上应用 Hölder 不等式得到
\[
 \left(|Q|^{-1}\int_Q\sum_{|k-j|<L}b_k^2\right)^{1/2}
 \le(2L)^{1/2-1/p}F(z),\qquad
 F=\left(\sum_k[M(b_k^2)]^{p/2}\right)^{1/p}.
\]
由于 $p/2>1$，由标量极大定理和 $B_p$ 在伸缩下的不变性可得
\[
 \|F\|_p^p\lesssim\sum_k\|b_k\|_p^p
 \le B_p^p\sum_k\|f_k\|_p^p\le B_p^p\|W\|_p^p.
\]
现在利用 Hilbert 平方范数的三角不等式，并在 $Q$ 上用 Cauchy--Schwarz 不等式，可得 $|h-c_Q|$ 的平均值不超过
\[
 C\left(C_*\delta MW(z)+B_2M(W^2)(z)^{1/2}
       +L^{1/2-1/p}F(z)\right).
\]
将 $c_Q$ 换成平均值 $h_Q$，至多增加一个二倍因子。因此，对二进立方体取上确界所定义的二进尖锐函数 $h^\#(z)=\sup_{Q\ni z}|Q|^{-1}\int_Q|h-h_Q|$，其 $L^p$ 范数由所述右端控制。此处用到极大算子在 $p$ 和 $p/2$ 上的有界性。应用尖锐函数不等式 $\|h\|_p\lesssim_p\|h^\#\|_p$，见 \cite{r-74}, Chapter~IV, Section~3.6, Corollary~1。其中对较小指数的要求成立，因为
\[
 \|h\|_2^2=\sum_k\|S^{(k)}f_k\|_{L^2(H)}^2
 \le B_2^2\sum_k\|f_k\|_2^2<\infty
\]
且 $2<p$。
